% SPDX-License-Identifier: Apache-2.0
% covers: AxiHost
\datasheet{AxiHost}{The host side of an AXI4 link: bursts out, identifiers handed out, answers back.}

\dsfacts{\code{lib/parts/src/bus/axi.rs}}{\code{txhdl\_parts::bus::axi::AxiHost}}%
{\code{//docs:axi}, \code{//docs:noc}}

\subsection*{Function}

An AXI4 link in \code{txhdl\_parts} has a tracker at each end, between
a client and the five AXI channels. \code{AxiHost} is the tracker on
the host side. A host client gives it an \code{Issue} on \code{issue}:
an address phase with no identifier, and a bit that says whether it is
a read. \code{AxiHost} gives the burst an identifier, sends it on
\code{ar} or \code{aw}, and tells the client the identifier on
\code{grant}.

It also passes the client's write beats from \code{wbeat} to \code{w},
turns each \code{B} beat into a \code{Done} on \code{done}, and passes
each read beat from \code{r} to \code{rdata}. The client hands an
identifier back on \code{release} when it has taken that burst's
answer. \code{axi}, \code{axi\_units} and \code{axi\_to\_unit} make the
channels it runs on. Its ports are named types, \code{HostIn} and
\code{HostOut}, so a design may supply a unit of its own instead.

\subsection*{Parameters}

\begin{center}\footnotesize
\begin{tabular}{@{}lp{0.7\textwidth}@{}}
\toprule
Parameter & Meaning \\
\midrule
\code{A} & The address width, in bits. \\
\code{D} & The data width, in bits. \\
\code{S} & The write strobe width, one bit per byte lane, so \code{D / 8}. \\
\code{I} & The identifier width, in bits. \\
\code{NIDS} & How many identifiers it hands out, so \code{1 << I}. It is
the width of \code{busy}, and the most bursts in flight at once. \\
\bottomrule
\end{tabular}
\end{center}

\code{S} and \code{NIDS} are stated rather than computed, because an
expression in a const parameter's position needs nightly Rust. The
source states that a link whose \code{S} and \code{NIDS} disagree with
its \code{D} and \code{I} will not behave. The tables below use
\code{AxiHost<32, 32, 4, 2, 4>}, the tracker of Vreteno's link and of
its board.

\dstables{AxiHost}

\subsection*{Behaviour}

\begin{itemize}
\item At each rising edge it looks at the issue at the head of
\code{issue}. The identifier that burst would take is \code{turn}.
\item The burst goes out when that identifier's bit of \code{busy} is
clear, \code{grant} has room, and the address channel it needs has
room: \code{ar} for a read, \code{aw} for a write.
\item A burst that goes out is sent with \code{turn} as its \code{id}.
Its \code{addr}, \code{len}, \code{size}, \code{burst}, \code{lock},
\code{cache}, \code{prot}, \code{qos} and \code{region} are the
issue's.
\item In that cycle \code{grant} sends the identifier, \code{turn}
moves on by one, and the identifier's bit of \code{busy} is set. At
most one burst goes out per cycle.
\item The identifiers go round robin. A burst waits for its own turn's
identifier to be free, and for nothing else. Another free identifier
does not let it go.
\item \code{release} is taken whenever it offers, and it clears the
bit of the identifier it names. An identifier is therefore outstanding
from its address phase until the client has taken its answer.
\code{//docs:axi} explains why freeing it on the response would be
wrong.
\item A write beat passes from \code{wbeat} to \code{w} unchanged when
\code{w} has room. AXI4 puts no identifier on \code{w}, so the source
says a client sends a burst's beats before it issues the next burst.
\item A \code{B} beat becomes a \code{Done} of the same \code{id} and
\code{resp} when \code{done} has room. A read beat passes from
\code{r} to \code{rdata} unchanged when \code{rdata} has room.
\end{itemize}

\subsection*{Verification}

\begin{itemize}
\item Four tests in \code{bus::axi} run \code{AxiHost<16, 32, 4, 2, 4>}
and \code{AxiPer<16, 32, 4, 2>} with a host client and a peripheral
client on the two ends:
\begin{itemize}
\item \code{a\_read\_reads\_what\_the\_write\_wrote};
\item \code{answers\_come\_back\_out\_of\_order};
\item \code{a\_transaction\_is\_answered\_by\_another\_process};
\item \code{every\_burst\_is\_answered\_and\_reads\_agree\_with\_writes},
which issues 24 bursts of one to four words to a peripheral that
\code{serve} runs with four slots.
\end{itemize}
\item The tests of \code{bus::axi::sim} put the same two trackers in
front of the simulated RAM:
\begin{itemize}
\item \code{a\_burst\_reads\_what\_a\_burst\_wrote};
\item \code{a\_strobe\_covers\_the\_lanes\_it\_says};
\item \code{a\_fixed\_burst\_stays\_on\_one\_word};
\item \code{an\_address\_past\_the\_end\_is\_answered};
\item \code{an\_image\_is\_loaded\_and\_read\_back}.
\end{itemize}
All of these run in \code{//lib/parts:parts\_test}.
\item \code{ex\_axi} issues three bursts and has them answered out of
order. The build simulates the netlist of \code{AxiHost<16, 32, 4, 2,
4>} against that run under nvc and Verilator:
\begin{itemize}
\item \code{//docs:axi\_sim\_axi\_host\_axi\_host\_tb\_test};
\item \code{//docs:axi\_vsim\_axi\_host\_test}.
\end{itemize}
\item The tests of the router, the AXI-Lite bridge, the Wishbone bridge
and the network on chip in \code{txhdl\_parts} run it in simulation.
So do \code{//ddr3:ddr3\_test}, \code{//gpu/razboj:razboj\_test},
\code{//soc:demo\_test} and \code{//cpu/vreteno:lockstep}.
\end{itemize}

\subsection*{Used by}

The examples \code{ex\_axi}, \code{ex\_axi\_serve}, \code{ex\_axi\_router},
\code{ex\_axi\_lite}, \code{ex\_axi\_reg}, \code{ex\_axi\_wb},
\code{ex\_eth} and \code{ex\_hdmi}. Vreteno's run in
\code{cpu/vreteno/src/run.rs} and \code{cpu/vreteno/src/main.rs}, and
the board's design, \code{Board}, as its field \code{host}. The
four-corner system in \code{soc/src/lib.rs}, once for the core and
once for Razboj. Razboj's simulation in \code{gpu/razboj/src/sim.rs}.

\subsection*{Limits}

At most \code{NIDS} bursts are in flight at once. A client that never
releases an identifier stops the unit when the turn reaches that
identifier. It sends \code{lock}, \code{cache}, \code{qos} and
\code{region} as the client gave them and acts on none of them; nothing
in \code{bus::axi} makes an exclusive access. It does not check
\code{S} and \code{NIDS} against \code{D} and \code{I}.
